\section{The Ecological Significance of Southeastern Australian Forests}

The forests of southeastern Australia span a remarkable climatic and edaphic
gradient, from the cool temperate rainforests of Tasmania and the Otway Ranges
to the dry sclerophyll woodlands of the western slopes of the Great Dividing
Range~\cite{keith2014}. These ecosystems support over 1{,}500 vascular plant
species and provide critical habitat for threatened fauna including the
Greater Glider (\textit{Petauroides volans}), the Powerful Owl
(\textit{Ninox strenua}), and the Long-footed Potoroo
(\textit{Potorous longipes})~\cite{bradshaw2012}.

Despite their ecological and economic importance—timber, water catchment,
carbon storage, and recreation alone contribute an estimated AUD~2.3~billion
annually to regional economies—these forests face unprecedented threats.
The 2019--2020 Black Summer bushfires burnt approximately 5.8~million hectares
of forest in New South Wales and Victoria, and projections under the
Representative Concentration Pathway~8.5 (RCP8.5) suggest a three- to
four-fold increase in the number of extreme fire-weather days by 2070
\cite{keith2014}. Concurrent pressures from invasive species, habitat
fragmentation, and changing hydrological regimes compound the challenge.

Effective conservation and management in this dynamic setting require reliable,
repeatable monitoring that can detect ecologically meaningful change against a
background of high natural variability. Yet the statistical methods available
to practitioners have not kept pace with the volume and complexity of data now
being collected through remote sensing, autonomous acoustic recorders, camera
traps, and environmental DNA (eDNA) metabarcoding.
